\section{Metodyka oceny złożoności}

Z uwagi na brak możliwości sprowadzenie tego procesu do pomiaru listy obiektywnych metryk, przedstawiona ocena złożoności ma charakter subiektywnego opisu wad i zalet danego rozwiązania.

\section{Metodyka pomiaru wydajności}

\subsection{Eksperyment}

Jako eksperyment definiuje się zbiór parametrów:

\begin{itemize}[leftmargin=*]
    \item Czy oraz jaki protokół jest badany
    \item Limit przepustowości
    \item Emulowana wartość delay oraz jitter
    \item Taktowanie rdzeni przypisanych do klienta i serwera
    \item Ilość rdzeni przypisanych do klienta i serwera
\end{itemize}

Cały test składa się z wielu pojedynczych eksperymentów.

\subsection{Metryki}

Pomiarem objęta została \textbf{przepustowość} oraz \textbf{wykorzystanie procesora}.

\subsection{Narzędzia}

Skrypt automatyzujący proces testowania został napisany z wykorzystaniem języka \textbf{Python}.\\

Proces budowy obrazu dysku został zautomatyzowany w oparciu o język \textbf{Bash}.\\

Pozostałe narzędzia:

\begin{enumerate}[leftmargin=*]
    \item \textbf{iperf3} - Narzędzie do pomiaru przepustowości sieci pomiędzy dwoma urządzeniami (architektura klient-serwer).
    \item \textbf{tc} - Zestaw narzędzi do sterowania przepływem ruchu sieciowego, pozwalający na symulację warunków sieciowych (np. opóźnienia, ograniczenia przepustowości).
    \item \textbf{KVM/QEMU} - Zintegrowany z jądrem Linux hipernadzorca (KVM) oraz emulator sprzętu (QEMU), umożliwiające pełną wirtualizację ze wsparciem sprzętowym.
    \item \textbf{virsh} - CLI do zarządzania maszynami wirtualnymi.
    \item \textbf{live-build} - Narzędzie do zautomatyzowanego budowania niestandardowych, uruchamialnych z nośników wymiennych obrazów systemów operacyjnych bazujących na Debianie.
    \item \textbf{cpupower} - Narzędzie do monitorowania i modyfikowania parametrów pracy procesora, w tym profili zarządzania energią i taktowaniem.
    \item \textbf{sar} - Narzędzie użyte do gromadzenia danych o obciążenie CPU.
\end{enumerate}

\newpage

\subsection{Założenia narzędzia testowego}

Podstawowym założeniem projektu jest wykorzystanie frameworku live-build do automatycznego przygotowania już skonfigurowanego systemu zawierającego gotowy do użycia skrypt oraz środowisko testowe.\\

Wszelka konfiguracja zgodna jest z praktyką infrastructure as code.\\

Całość maksymalizuje potencjał do otworzenia przez całkowitą eliminację manualnej ingerencji poprzez automatyzację nie tylko procesu testowania jak i również samego etap budowania narzędzia.\\

Sercem projektu jest skrypt orkiestrujący. Jego zadaniem jest przeprowadzenie testu poprzez wywoływanie serii poleceń systemowych na maszynach wirtualnych oraz hoście.

\subsection{Infrastruktura}

Test zakłada wykorzystanie trzech maszyn wirtualnych:\\

Klient \rightleftarrows \ Router \rightleftarrows \ Serwer\\

Każda z maszyn uruchamiana jest z użyciem identycznego obrazu dysku, przygotowywanego w analogiczny sposób co obraz hosta.\\

To jaką rolę przyjmie dana maszyna jest więc umowne, zależne od uruchamianych komend i konfiguracji maszyn wirtualnych z poziomu hosta.\\

Klient i serwer połączeni są odrębnymi sieciami z routerem.\\
Przekierowywanie ruchu pomiędzy sieciami odbywa się całkowicie na routerze.\\

Istnieje również trzecia, odrębna sieć, dedykowana zarządzaniu maszynami, która łączy wszystkie maszyny wirtualne wraz z hostem.